\section{Introduction}


The construction industry contributes more than 10\% of the GDP of many countries \citep{RILEM2022GLOBEEnvironment}, along with 20\% of global CO\textsubscript{2} emissions \citep{Huang2018CarbonSector}.
From both economic and environmental viewpoints, there is therefore much to gain by optimising structural designs in terms of minimal costs and emissions, while ensuring they are safe enough from a societal point of view.
Decision makers in the field of structural engineering, such as designers, code committees, or asset managers, benefit from guidance on this optimisation.
As such, a crucial question from their perspective is the appropriate level of reliability for a structure, in other words, \emph{how safe should a structure be?}
The answer to this question then informs decisions on structural design.
To find it, the costs, risks, and benefits of a structure need to be considered, ideally in the broad sense, including environmental impacts and societal risks due to the structure potentially failing and causing harm to its users.
These aspects can be quantified through life cycle analyses, taking all three `pillars' of sustainability into account \citep{Knibbe2026PlaceHolder, Valdivia2021PrinciplesAssessment}.
The level of reliability corresponding to the best `performance' over the life cycle of the structure can then be considered as an appropriate design target for decision makers.

This level can be numerically represented by the target failure probability $P_{f,target}$ or, equivalently, the target reliability index $\beta_{target}$.
These are linked through the following conversion:
\begin{equation}
\label{eqn:beta_Pf}
    \beta = \Phi^{-1}(1-P_f)
\end{equation}
Where $\Phi^{-1}(\cdot)$ is the inverse of the cumulative distribution function of a standard-normal distributed random variable.
Current norms on target reliabilities \citep{iso2394, EN1990} distinguish optimal and acceptable reliability values \citep{Fischer2019OptimalDesign}, noted here as $\beta_{opt}$ and $\beta_{acc}$ respectively, where the most stringent value forms the target reliability:
\begin{equation}
    \beta_{target} = \max(\beta_{opt}, \beta_{acc})
\end{equation}
The optimal value is found by minimising a lifetime cost function, and the acceptable value through the minimal required investment to ensure adequate life safety.
This current framework has been developed based on the work of several authors in the field of structural safety, including \citep{Rackwitz2000OptimizationVerification, Rackwitz2006TheCriteria, Rackwitz2008TheVerification, Viljoen2012DerivingLQI, Fischer2019OptimalDesign, Pandey2004LifeSafety, Pandey2024LifeStructures}.
At present, however, the cost functions used to determine both $\beta_{opt}$ and $\beta_{acc}$ are based on monetary considerations only.
Environmental concerns are not accounted for, which is undesirable considering the aforementioned impact of construction on the environment and the societal urgency of mitigating climate change.
If the construction sector is to contribute to this mitigation, decision-makers require an approach for determining target reliabilities that \emph{does} take environmental impacts into account.
Such an approach is presented in this paper, focusing on optimal reliabilities.
This builds upon previous work by the same authors, focused on acceptable reliability \citep{Knibbe2026PlaceHolder}. 

Sections 2 and 3 present the current approach for finding optimal reliability values, and expand on the need for further development from social and environmental perspectives.
This serves as background to Section 4, where the augmented approach itself is presented, followed by a demonstration through a typical case study in Section 5.
Section 6 contains an analysis of the case study results, and the impact of the presented approach with respect to the conventional approach.
Finally, Section 7 provides a conclusion and a brief discussion.
